\documentclass[11pt,a4paper]{article}

\usepackage{amsmath,amssymb}
\usepackage{booktabs}
\usepackage[colorlinks=true,linkcolor=blue,citecolor=blue,urlcolor=blue]{hyperref}
\usepackage{geometry}
\geometry{margin=2.5cm}
\usepackage{microtype}

\newcommand{\Es}{E^{*}}
\newcommand{\Order}{\mathcal{O}}
\newcommand{\norm}[1]{\left\lVert #1 \right\rVert}
\newcommand{\abs}[1]{\left\lvert #1 \right\rvert}

\title{A matrix-free $\mathcal{H}^2$/FMM operator for the elastic
half-space flexibility problem\\[2pt]
\large strategy, design choices, and validation}
\author{ASPHER}
\date{2026-06-27 --- revised 2026-07-13 (active-set solve)}

\begin{document}
\maketitle

\begin{abstract}
This note documents the matrix-free hierarchical operator
(\texttt{backend="h2"}) used by ASPHER to apply the Boussinesq half-space
flexibility matrix without ever forming it. The operator is a black-box
fast multipole method (bbFMM) built on tensor-product Chebyshev
interpolation. We explain why this design is preferred over a dense matrix
or a classical $\mathcal{H}$-matrix with ACA-compressed blocks, describe the
algorithm and the choices that make it efficient for a regular grid with a
translation-invariant kernel, and report measured memory, accuracy, and a
comparison against the FFT solver Tamaas. A final section documents the
masked matvec and the active-set solve built on it (July 2026), which cut the
large-grid solve by a further $\sim\!4\times$ in time and $\sim\!2\times$ in
memory without changing the solution.
\end{abstract}

\section{Problem and discretisation}

We solve normal contact of a rigid rough indenter against an elastic
half-space. The surface displacement under a pressure field is the
convolution
\begin{equation}
  u_z(\mathbf{x}) = \int G(\mathbf{x}-\mathbf{x}')\,p(\mathbf{x}')\,\mathrm{d}A',
  \qquad
  G(r) = \frac{1}{\pi \Es\, r}.
\end{equation}
On a uniform $N_s\times N_s$ grid of square elements of side $h=L/N_s$,
discretising with piecewise-constant pressure gives the influence
matrix $S_{ij}$ equal to the Love~\cite{Love1929} closed-form integral of $G$
over a source element, evaluated at a target element centre. Because the geometry
is a regular grid, $S_{ij}$ depends only on the index offset
$(|i_x-j_x|,|i_y-j_y|)$ -- the operator is \emph{translation invariant}, so
all $N^2$ entries ($N=N_s^2$) are served from one $N_s\times N_s$ lookup
table built with $\Order(N)$ kernel evaluations.

The contact problem is the bound- and equality-constrained quadratic
program $\min \tfrac12 \mathbf{p}^\top S\,\mathbf{p} + \mathbf{p}^\top
\mathbf{g}_0$ s.t.\ $\mathbf{p}\ge 0$, $\mathrm{mean}(\mathbf{p})=\bar p$,
solved by the Polonsky--Keer projected conjugate gradient~\cite{PolKeer1999}
(see \texttt{doc/theory/pcg.tex}). \textbf{Each PCG iteration needs two
applications of $S$.} The operator is therefore used only as a
matrix--vector product; we never need explicit matrix entries, an LU
factorisation, or a stored matrix. This is the key observation that selects
the algorithm.

\section{Why matrix-free: the cost landscape}

\begin{table}[h]
\centering
\begin{tabular}{lll}
\toprule
approach & memory & matvec \\
\midrule
dense matrix                       & $\Order(N^2)$        & $\Order(N^2)$ \\
classical $\mathcal{H}$-matrix (ACA) & $\Order(N\log N)$  & $\Order(N\log N)$ \\
FFT convolution (zero-padded)      & $\Order(N)$          & $\Order(N\log N)$ \\
$\mathcal{H}^2$/FMM (this note)    & $\Order(N)$          & $\Order(N)$ \\
\bottomrule
\end{tabular}
\caption{Asymptotic cost. For $N_s=1024$ ($N\approx10^6$) the dense matrix
alone is $\approx 8.8$\,TB, so it is impossible; the classical
$\mathcal{H}$-matrix needs several GB.}
\end{table}

The previous ASPHER backend was a classical $\mathcal{H}$-matrix~\cite{Hackbusch2015}:
a list of admissible blocks, each compressed independently as $A_{ts}\approx U_{ts}
V_{ts}^\top$ by adaptive cross approximation (ACA). This works, but for a
regular half-space kernel it stores \emph{redundant bases}: many blocks at
the same tree level with the same relative offset are the \emph{same matrix},
yet each is approximated and stored separately. The memory is
$\Order(N\log N)$ with a large constant (e.g.\ $6.2$\,GB at $N_s=512$).

The $\mathcal{H}^2$/FMM design removes this redundancy in two ways:
(i)~it stores a \emph{single basis per cluster} rather than factors per
block, replacing $A_{ts}\approx U_{ts}V_{ts}^\top$ with
$A_{ts}\approx V_t\, S_{ts}\, V_s^\top$ where $V_t,V_s$ are shared
interpolation operators and $S_{ts}$ is a small coupling matrix; and
(ii)~it exploits translation invariance so that the $V$'s and $S_{ts}$'s are
\emph{cached by $(\text{level},\text{relative offset})$}. The result is
$\Order(N)$ memory with a small constant.

\section{The strategy: black-box FMM with Chebyshev interpolation}

\subsection{Uniform quad-tree}
A balanced quad-tree partitions the grid down to square leaves of side
$\ell$ (\texttt{leaf\_side}, default~$8$, i.e.\ $64$ DOF/leaf). Boxes store
only \emph{index ranges} $[i_{x0},i_{x0}+s)\times[i_{y0},i_{y0}+s)$, never
explicit index lists, so the tree is $\Order(N/\ell^2)$ tiny structs.

\subsection{Near/far split}
A target box interacts \emph{directly} (exact kernel) with leaves in its
\texttt{near\_radius} neighbourhood (default $1$, the $3\times3$ block) and
\emph{approximately} (interpolation) with the standard FMM interaction list
-- children of the parent's neighbours that are not themselves near
neighbours. The interaction list has bounded size at every level, and
together the near field and the per-level interaction lists cover every pair
$(i,j)$ exactly once.

\subsection{Degenerate-kernel (Chebyshev) approximation}
For a well-separated target box $t$ and source box $s$, the kernel is
approximated by interpolating each argument on $q$ Chebyshev nodes
$\{\xi_a\}$ per axis:
\begin{equation}
  G(\mathbf{x},\mathbf{y}) \approx
  \sum_{a}\sum_{b} S_q(\mathbf{x},\boldsymbol{\xi}_a^{t})\,
  G(\boldsymbol{\xi}_a^{t},\boldsymbol{\xi}_b^{s})\,
  S_q(\mathbf{y},\boldsymbol{\xi}_b^{s}),
\end{equation}
with the 1-D interpolation weight
$S_q(x,\xi_m)=\tfrac1q+\tfrac2q\sum_{n=1}^{q-1}T_n(\xi_m)T_n(x)$
($T_n$ Chebyshev polynomials). In 2-D the operators are tensor products
($r=q^2$ nodes per box), applied as two 1-D passes rather than a Kronecker
product. This is exactly the black-box FMM~\cite{FongDarve2009}: only kernel
\emph{evaluations} are needed, never analytic multipole expansions.

\subsection{The six passes}
A matvec $\mathbf{u}=S\,\mathbf{p}$ runs
\[
\text{P2M}\;\to\;\text{M2M}\!\uparrow\;\to\;\text{M2L}\;\to\;
\text{L2L}\!\downarrow\;\to\;\text{L2P}\;\;+\;\;\text{near}.
\]
\textbf{P2M} compresses leaf source values onto Chebyshev coefficients
$M_s = W^\top X_s W$. \textbf{M2M} aggregates children to parents,
$M_{\text{par}}\mathrel{+}=\sum_c R_c M_{\text{child}}$. \textbf{M2L}
translates source to target coefficients along the interaction list,
$L_t\mathrel{+}=\sum_s K_{ts}M_s$. \textbf{L2L} propagates parents to
children, $L_{\text{child}}\mathrel{+}=R_c^\top L_{\text{par}}$. \textbf{L2P}
evaluates $\mathbf{u}_t\mathrel{+}=W L_t W^\top$ at leaf DOFs. The
\textbf{near} pass adds direct contributions of neighbouring leaves. All
passes are OpenMP-parallel over boxes/leaves with disjoint writes.

\subsection{Translation-invariant caching (the decisive choice)}
On a uniform grid with a translation-invariant kernel:
\begin{itemize}
\item the \textbf{M2M/L2L} transfer matrices depend only on the child
quadrant and $q$ -- there are exactly \textbf{four} $q^2\times q^2$ matrices,
reused at every level (boxes are normalised to $[-1,1]^2$ by their geometric
extent, giving exact $1/2$ scaling between levels);
\item the \textbf{M2L} coupling $K_{ts}[a,b]=g(\boldsymbol{\xi}_a^t-
\boldsymbol{\xi}_b^s)$ depends only on $(\text{level},\,\Delta b_x,\,\Delta
b_y)$ and is cached -- a few hundred distinct $q^2\times q^2$ matrices total;
\item the \textbf{near} stencils depend only on the relative leaf offset --
at most $(2\,\text{near\_radius}+1)^2$ distinct $\ell^2\times\ell^2$ dense
blocks.
\end{itemize}
The far-field point kernel is taken to be the \emph{same} Love cell integral
as the near field, $g(\mathbf{d})=\mathrm{love\_uz}(\mathbf{d},h/2,h/2)/(\pi
\Es)$, evaluated at continuous node offsets. Using the Love integral (rather
than the bare $1/r$) makes near and far consistent, so the only far-field
error is the interpolation error.

\section{Complexity and measured memory}

The caches are $\Order(\log N)$ (couplings) and $\Order(1)$ (near stencils);
the per-box coefficient buffers are $\Order(N/\ell^2)$. Total memory is
$\Order(N)$, dominated asymptotically by the buffers at $\approx 13$
bytes/DOF. Table~\ref{tab:mem} (\texttt{bench\_h2\_memory.py}, $q=6$) shows
the bytes/DOF \emph{decreasing} to this bound while a dense matrix explodes.

\begin{table}[h]
\centering
\begin{tabular}{rrrrrr}
\toprule
$N_s$ & $N$ & coupling & buffers & total & dense $N^2$ \\
      &     & [MiB]    & [MiB]   & [MiB] & [MiB] \\
\midrule
64   & 4\,096      & 0.79 & 0.05  & 1.12  & 128 \\
128  & 16\,384     & 1.19 & 0.19  & 1.66  & 2\,048 \\
256  & 65\,536     & 1.58 & 0.75  & 2.61  & 32\,768 \\
512  & 262\,144    & 1.98 & 3.00  & 5.26  & 524\,288 \\
1024 & 1\,048\,576 & 2.37 & 12.0  & 14.65 & 8\,388\,608 \\
2048 & 4\,194\,304 & 2.77 & 48.0  & 51.05 & 134\,217\,728 \\
\bottomrule
\end{tabular}
\caption{H2 operator memory ($q=6$, $\ell=8$, near\_radius~$=1$). The near
stencil cache is a constant $0.28$\,MiB. At $N_s=2048$ the operator is
$51$\,MiB versus $128$\,TiB for the dense matrix.}
\label{tab:mem}
\end{table}

\section{CPU-time scaling to $N_s=16384$}

Table~\ref{tab:cputime} (\texttt{bench\_h2\_cputime.py}, $q=6$, 20-core node)
reports the operator build time, per-matvec wall time, and peak process RSS
for the \texttt{example\_rough\_contact} problem from $N_s=128$ up to
$N_s=16384$ ($N\approx2.7\times10^8$ unknowns), with the full rough-surface
solve included where a converged solution is practical.

\begin{table}[h]
\centering
\begin{tabular}{rrrrrrr}
\toprule
$N_s$ & $N$ & build [s] & matvec [s] & RSS [GiB] & solve [s] & iters \\
\midrule
128   & 16\,384       & 0.007 & 0.0005 & 0.04 & 0.21  & 90 \\
256   & 65\,536       & 0.011 & 0.0013 & 0.05 & 0.80  & 199 \\
512   & 262\,144      & 0.023 & 0.0052 & 0.08 & 3.64  & 241 \\
1024  & 1\,048\,576   & 0.064 & 0.019  & 0.19 & 48.0  & 906 \\
2048  & 4\,194\,304   & 0.245 & 0.098  & 0.26 & ---   & --- \\
4096  & 16\,777\,216  & 0.92  & 0.34   & 0.92 & ---   & --- \\
8192  & 67\,108\,864  & 3.65  & 1.39   & 3.57 & ---   & --- \\
16384 & 268\,435\,456 & 13.98 & 6.06   & 14.2 & ---   & --- \\
\bottomrule
\end{tabular}
\caption{H2 CPU-time study ($q=6$, all cores). Build and per-matvec both scale
$\approx\Order(N)$ (each $4\times$ in $N$ costs $\approx4\times$). The
$2.7\times10^8$-DOF operator builds in $14$\,s, applies in $6$\,s, and fits in
$14$\,GiB; the dense matrix would need $\sim5\times10^{8}$\,GiB. Full solves
are listed to $N_s=1024$: beyond that the projected-CG iteration count grows,
so the total solve time -- not the operator -- becomes the limiting cost, and
only the operator metrics are reported.}
\label{tab:cputime}
\end{table}

Build and matvec scale linearly in $N$, confirming the $\Order(N)$ design; the
peak RSS is $\approx57$ bytes/DOF (the intrinsic caches are $\approx12$
bytes/DOF, the remainder being interaction-list indices and the interpreter).
The operator therefore reaches grids that are entirely out of reach for a
dense or classical $\mathcal{H}$-matrix representation, at second-scale build
and matvec times.

\section{Accuracy}

The interpolation error decreases geometrically with $q$; the near field is
exact. Against the dense Love operator ($N_s=32$, \texttt{tests/test\_h2.cpp}):
relative $L_2$ matvec error is $1.3\times10^{-4}$ at $q=4$ and
$3.2\times10^{-6}$ at $q=6$. The H2 backend, plugged into the same PCG,
reproduces the Hertz contact radius and peak pressure and matches the
$\mathcal{H}$-matrix contact area exactly.

\section{Comparison with the FFT solver (Tamaas)}

For a fully active rectangular grid the operator is a convolution and an FFT
matvec ($\Order(N\log N)$, $\Order(N)$ memory) is the natural reference;
Tamaas implements this via a zero-padded \texttt{dcfft} operator.
Table~\ref{tab:tamaas} (\texttt{compare\_tamaas\_h2.py}) compares the H2
backend ($q=6$) against Tamaas~2.8.1 on the same self-affine surfaces.

\begin{table}[h]
\centering
\begin{tabular}{rrrrr}
\toprule
$N_s$ & frac.\ (Tamaas) & frac.\ (H2) & rel.\ $L_2$ pressure & H2 mem [MiB] \\
\midrule
64  & 0.1265 & 0.1260 & 3.26\% & 1.12 \\
128 & 0.1215 & 0.1201 & 3.17\% & 1.66 \\
256 & 0.1112 & 0.1098 & 2.92\% & 2.61 \\
\bottomrule
\end{tabular}
\caption{H2 vs Tamaas FFT. Contact fractions agree to $\sim10^{-3}$.}
\label{tab:tamaas}
\end{table}

Timing (single thread, \texttt{OMP\_NUM\_THREADS=1}) is reported in
Table~\ref{tab:tamaas-time}. The H2 per-matvec cost is already below the FFT
at every size, and the ratio improves with $N$ (from $0.78$ at $N_s=64$ to
$0.69$ at $N_s=512$) as the $\Order(N)$ apply overtakes the FFT's
$\Order(N\log N)$. End-to-end the H2 solve is $\sim1.5\times$ faster here,
which also reflects a (somewhat) lower iteration count of the projected CG.

\begin{table}[h]
\centering
\begin{tabular}{rrrrrrr}
\toprule
& \multicolumn{3}{c}{per-matvec [ms]} & \multicolumn{3}{c}{end-to-end solve [s]} \\
\cmidrule(lr){2-4}\cmidrule(lr){5-7}
$N_s$ & Tamaas & H2 & H2/T & Tamaas & H2 & H2/T \\
\midrule
64  &  0.36 &  0.28 & 0.78 &  0.029 & 0.017 & 0.58 \\
128 &  1.63 &  1.29 & 0.79 &  0.204 & 0.128 & 0.63 \\
256 &  6.92 &  5.59 & 0.81 &  1.230 & 0.870 & 0.71 \\
512 & 37.30 & 25.89 & 0.69 & 10.319 & 6.822 & 0.66 \\
\bottomrule
\end{tabular}
\caption{Single-thread timing, H2 ($q=6$) vs Tamaas dcfft FFT
(\texttt{compare\_tamaas\_h2.py}). Per-matvec is the same operation
$\mathbf{u}=S\mathbf{p}$ on both sides.}
\label{tab:tamaas-time}
\end{table}

The residual $\sim3\%$ pressure $L_2$ difference is \emph{not} an H2 error:
the H2 operator reproduces the exact Love operator to $3\times10^{-6}$
($q=6$), whereas Tamaas~2.8.1's \texttt{dcfft} influence coefficients
oscillate $\pm2$--$8\%$ from the exact Love values at $1$--$3$ cell
separations (Gibbs-like), and its effective modulus is $2\Es{}^2$ (both
verified against Hertz theory; see the project README). The difference is
therefore dominated by the FFT reference, and it is the same $\approx3.3\%$
seen for the $\mathcal{H}$-matrix backend. Both H2 and FFT are $\Order(N)$ in
memory; H2 is additionally \emph{local}, which allows skipping work outside
the contact set --- something a global FFT cannot do without re-padding.
This is no longer just ``room'': the masked matvec and the active-set solve
of Section~\ref{sec:active-note} realise it.

\section{Convergence acceleration}
\label{sec:accel}

The $\Order(N)$ operator makes a single matvec cheap, but the \emph{number} of
projected-CG iterations grows with $N_s$ (Table~\ref{tab:cputime}), so for large
grids the iteration count -- not the operator -- becomes the bottleneck. Two
complementary, optional accelerations address this; both leave the result
unchanged and \texttt{precond="none"} reproduces the original solver bit-for-bit.

\paragraph{Why iterations grow.}
The Boussinesq operator has Fourier symbol $\widehat S(\mathbf q)\propto
1/\abs{\mathbf q}$. On an $N_s\times N_s$ grid the resolved wavenumbers span
$q_{\min}\sim 2\pi/L$ to $q_{\max}\sim\pi N_s/L$, so the spectral condition number
is $\kappa(S)=\widehat S(q_{\min})/\widehat S(q_{\max})=q_{\max}/q_{\min}\sim N_s$
and unpreconditioned CG needs $\sim\sqrt{\kappa}\sim\sqrt{N_s}$ iterations --
intrinsic to the kernel, on top of active-set discovery. The smooth Hertz problem
grows only mildly with $N_s$, confirming the conditioning term; the very steep
$90\!\to\!906$ growth in Table~\ref{tab:cputime} was mostly the
\emph{surface roughening} with $N_s$ (the example surface adds higher-frequency
content as the grid refines). On a band-limited surface with a \emph{fixed
physical} spectrum the baseline grows $\approx\sqrt{N_s}$
(Table~\ref{tab:precond}, $26\!\to\!180$ over a $16\times$ range in $N_s$).

\paragraph{Spectral (Fourier) preconditioner.}
Choosing $M^{-1}$ with the inverse symbol $1/\widehat S(\mathbf q)\propto
\abs{\mathbf q}$ collapses $\kappa$ to $\Order(1)$. The bare wavenumber form
$\abs{\mathbf q}$ and the circulant form $1/\widehat S_c$ (the FFT of the kernel's
influence coefficients) coincide here because $\widehat S_c\sim 1/\abs{\mathbf q}$;
the overall scale of $M^{-1}$ is irrelevant (it cancels in CG), and the $q{=}0$
mode is zeroed (the mean/total load is fixed by the constraint, not by CG). Each
iteration applies $z=M^{-1}g$ by FFT to the residual \emph{masked to the contact
set}, keeps $z$ on contact, and removes its contact-set mean
(\texttt{fourier\_precond.cpp}, Eigen FFT). The change to the Polonsky--Keer PCG
is minimal: the preconditioned residual $z=M^{-1}g$ replaces $g$ in the direction
build, with $\beta=\max\!\big(0,\langle z,\,g-g_{\text{prev}}\rangle\big)/
\langle z_{\text{prev}},g_{\text{prev}}\rangle$ (PR$^+$ in the $M$-inner product)
and $t=z+\beta t$; the \textbf{exact line search}
$\tau=\langle g,t\rangle/\langle (St-\overline{St}),t\rangle$ and the projection,
overlap correction, and load balance are \emph{unchanged}.

\paragraph{Warm start and nested-grid (cascadic/FMG) continuation.}
The problem is solved coarse$\to$fine. The surface is \emph{restricted}
($2\times2$ block average) to every coarser level so each level is the
\emph{same physical} problem; the converged coarse pressure is \emph{prolonged by
injection} (replicate each value into its $2\times2$ block) to seed the next
finer level via \texttt{p\_init}, which carries both the pressure and the active
set. Injection is load-preserving and keeps a sharp contact boundary; bilinear
prolongation smears pressure across the contact edge and gives a worse active-set
seed. Coarse levels use a looser (cascadic) tolerance. The overlap correction can
still reactivate nodes on the fine level, so the warm start is a seed, not a
constraint. The whole coarse$\to$fine continuation is a single C++ entry point
with no Python orchestration: \texttt{hmc::solve\_contact\_nested(...)}
(\texttt{include/nested\_solve.hpp}), exposed to Python as
\texttt{hc.solve\_nested(grid\_size, gap, p\_nominal, coarsest=64, q=6, ...)};
it builds the grid hierarchy and per-level $\mathcal{H}^2$ operators internally,
restricts the gap by $2\times2$ averaging, and warm-starts each level by
injecting the prolonged coarse pressure (requiring
$\texttt{grid\_size}=\texttt{coarsest}\cdot2^k$), returning a standard
\texttt{ContactResult}.

\paragraph{Results.}
On a fixed-band rough surface ($p_{\text{nominal}}=0.05$, tol $10^{-8}$, $q=6$),
the preconditioner roughly halves the growth rate and the warm start cuts a
further $\sim1.4\times$ (Table~\ref{tab:precond}). The two combine to
$\sim4\times$ fewer fine iterations at $N_s=1024$ ($180\!\to\!45$), and the
\emph{entire} coarse$\to$fine solve costs less than a single cold solve on the
target grid (work ratio $0.79\times$ at $N_s\ge512$). The C++ port matches the
Python prototype and is solution-identical (contact-area change $0$, pressure
rel-$L_2\sim5\times10^{-7}$); at $N_s=1024$ the wall time drops from $10.2$\,s
(none) to $7.9$\,s (Fourier) to $6.5$\,s (nested), $\sim1.6\times$. For Hertz at
$N_s=64$ the preconditioner alone cuts $27\!\to\!16$ iterations, and warm-starting
from the converged solution converges in $0$ iterations.

\begin{table}[h]
\centering
\begin{tabular}{rrrrrrr}
\toprule
$N_s$ & baseline & $\abs{q}$ & circ.\ & speedup & nested (inj.) & nested (bilin.) \\
\midrule
64   & 26  & 19 & 19 & 1.37$\times$ & ---  & ---  \\
128  & 43  & 25 & 24 & 1.79$\times$ & ---  & ---  \\
256  & 56  & 33 & 33 & 1.70$\times$ & 26   & 30   \\
512  & 83  & 43 & 43 & 1.93$\times$ & 31   & 32   \\
1024 & 180 & 62 & 62 & 2.90$\times$ & 45   & 51   \\
\bottomrule
\end{tabular}
\caption{Projected-CG iteration counts on the fixed-band rough surface
(\texttt{experiments/precond\_pcg.py}, \texttt{nested\_grid.py}). ``baseline''
is the unpreconditioned solver; $\abs{q}$ and ``circ.'' are the spectral
preconditioner with the two symbols; ``nested'' is the fine-grid iteration count
of the cascadic continuation (coarsest $N_s{=}64$). Injection prolongation beats
bilinear at every size.}
\label{tab:precond}
\end{table}

\paragraph{Honest limitation.}
The method is \emph{not} fully grid-independent. Pure spectral preconditioning is
limited by the contact restriction: masking the residual to the contact set
breaks the FFT diagonalisation near the contact boundary and for fragmented
contact. And each octave of refinement resolves genuinely \emph{new} asperities,
so there is real new fine-scale work per level (fine iters $26\!\to\!45$). These
accelerations are the path toward near grid-independence; a true
monotone/constraint multigrid for the variational
inequality~\cite{Kornhuber1994,Brandt1977,BornemannDeuflhard1996} is the heavier
alternative, deferred for now.

\section{Active-set solve: masked matvec and restricted PCG}
\label{sec:active-note}

At light loads the contact set is a small fraction of the grid
($N_c/N\sim10^{-2}$--$10^{-3}$), yet the solver of
Section~\ref{sec:accel} still applies the operator and every $\Order(N)$
vector pass on the full grid. The active-set solve
(\texttt{hc.solve\_nested(..., active\_set=True)}, July 2026) exploits the
locality of the $\mathcal{H}^2$ apply to restrict the whole finest-level
iteration to a \emph{candidate set} and certify the result.

\paragraph{Masked matvec.} \texttt{H2Mask} is a flat per-box occupancy
bitmap over the quad-tree, OR-propagated to ancestors (parent-closed).
The masked apply runs P2M/M2M only on source-occupied boxes, skips
non-source entries inside the M2L interaction-list walk, and runs
L2L/L2P/near only on target boxes. For an input that vanishes off the
source leaves the output is \textbf{bit-for-bit identical} to the unmasked
matvec on the target leaves: every skipped contribution is an exact zero and
the kept summation order is unchanged (asserted in
\texttt{tests/test\_active.cpp}, double and float). Measured cost with source
$=$ target $=$ candidates: $1.8$--$2.4\%$ of a full matvec at
$N_c/N\approx10^{-2}$ ($N_s=2048/4096$) and $\approx\!0.9\%$ at $N_s=16384$;
the verification mode (full target grid) costs $2$--$10\%$ and runs once or
twice per solve, not per iteration.

\paragraph{Restricted Polonsky--Keer.} The finest nested level runs the
\emph{identical} PCG (Section~\ref{sec:accel}) with all $\Order(N)$ loops
driven by the candidate index list and all CG state compressed to
$\Order(N_c)$ slot-blocked vectors ($\approx2$--$4\,N_c$ entries). The
candidate set is $\mathcal{K}=\mathrm{dilate}(\text{prolonged coarse
contact})\cup\{\text{coarse gap}<\delta\}$; the gap term is mandatory
($\sim\!30\%$ of fine contact islands have no coarse-contact precursor on
rich spectra), and $\delta$ is kept generous because a tight threshold can
pass verification while boundary pressures are subtly wrong. Verification is
one streamed masked-source/full-target matvec plus a gap scan; violations
are dilated into $\mathcal{K}$ and the solve resumes warm-started (one round
suffices in practice), with a full-solve fallback as escape hatch. Two
findings shape the design: restriction does \emph{not} improve conditioning
(the islands span the domain, so the smooth inter-island modes keep their
domain-scale $\kappa$ --- the $\abs{q}$ preconditioner \textbf{stays on},
applied through its full-grid FFT with compressed scatter/gather), and the
reduced-wave-vector (capped-symbol $\min(\abs{q},q_c)$) preconditioner
variant of the FFT-homogenisation literature~\cite{GierdenEtAl2022} scales
like the unpreconditioned iteration and was rejected. Full theory, tables,
and the certification subtleties: \texttt{pcg.tex} \S``Active-Set
Acceleration'' and \texttt{h2\_fmm\_detailed.tex}.

\paragraph{Measured.} Same solution as the standard nested solve (pressure
rel-$L_2\sim10^{-14}$ double, $\Delta A_c=0$). At $N_s=4096$ (double):
$173\to46.5$\,s ($3.7\times$), solve-peak RSS $1593\to841$\,MB. At
$N_s=16384$ (float, full-band surface): $1445\to308$\,s ($4.7\times$),
$18.3\to10.9$\,GiB ($1.67\times$); in \emph{double}, the active-set solve
runs in $12.5$\,GiB where the standard double solve exceeds the machine's
31\,GiB entirely.

\section{Design choices, summarised}
\begin{itemize}
\item \textbf{bbFMM over analytic FMM}: kernel-agnostic, needs only
evaluations -- robust to the Love near-field form.
\item \textbf{Love far kernel, not $1/r$}: near/far consistency; error is
interpolation-only.
\item \textbf{Geometric box extent}: exact half-scaling so M2M/L2L collapse
to four cached matrices.
\item \textbf{Cache by (level, offset)}: turns $\Order(N\log N)$ stored
factors into $\Order(\log N)$ couplings $+\,\Order(1)$ stencils.
\item \textbf{Defaults} $q=6$, $\ell=8$, near\_radius~$=1$: $\sim10^{-6}$
accuracy at a few MiB; tunable per problem.
\item \textbf{Additive backend}: the dense and $\mathcal{H}$-matrix paths and
all existing tests are untouched; PCG plugs into \texttt{H2Operator::matvec}.
\end{itemize}

\section{Limitations and outlook}
Of the extensions once listed here, three are now implemented without
changing the public API: active-domain \emph{masking}
(Section~\ref{sec:active-note}); \emph{single-precision} execution
(\texttt{single\_precision=True} keeps float copies of the caches and runs
the matvec, PCG, and preconditioner FFT in \texttt{float} with all scalar
reductions accumulated in double --- see \texttt{pcg.tex}); and an exact
\emph{FFT-convolution backend} (\texttt{backend="fft"}, zero-padded Love
kernel; matches the dense operator to $\sim10^{-15}$, roughly at parity with
H2 in matvec time at $N_s=4096$, with H2 preferred at very large $N_s$ for
its $\Order(N)$ cost and smaller working set). Remaining next steps:
rectangular ($n_x\neq n_y$) and locally refined grids; leaf-side/$q$
auto-tuning; GPU offload of the near field and M2L.

\bibliographystyle{unsrt}
\bibliography{references}

\end{document}
